Cada proyecto deja una lección para el PFC. En PresustentacionesAWPS vimos que un mismo servicio de negocio puede tener dos contratos sin duplicar la lógica, y que los códigos HTTP importan: el REST devolvía \texttt{200} con datos vacíos para una solicitud que no existía y lo cambiamos a \texttt{404}. En SGIRI, que todo servicio externo necesita un tiempo límite y un plan B. En la gestión documental, que los métodos idempotentes permiten reintentar sin duplicar documentos. Para cerrar: SOAP y REST no compiten en todo, el contexto decide; y como ningún servicio web oculta del todo sus fallos, el cliente siempre debe estar preparado.
